\section{Self-Attributing}
\label{sec:self-attribution}

In this section, we introduce the final stage of \texttt{AgentEvolver}, \texttt{self-attributing} mechanism.
Existing methods such as GRPO~\citep{deepseek_math} rely on sparse, trajectory-level rewards that equally credit all actions, failing to distinguish between critical decisions and inconsequential steps. This  credit assignment leads to inefficient sample utilization and limits the policy's ability to learn from fine-grained feedback.

To address this challenge, \texttt{self-attributing} transforms multi-step trajectories into efficient learning signals by leveraging an LLM's reasoning capabilities to assess the contribution of each action. It serves as the concrete implementation of our proxy reward function, $F_\mathrm{reward}$ in Eq.~\ref{eq:reward_mapping} (Section~\ref{sec:formulation}). We reformulate credit assignment by shifting from temporal signal propagation to \emph{contribution-based attribution}. It leverages a large language model's reasoning capabilities to retrospectively determine whether each action contributed positively or negatively to the final outcome. 

We structure the learning signal around two distinct dimensions: \textbf{process quality} (the contribution of an action step) and \textbf{outcome effectiveness} (the final trajectory success). To preserve the distinct contribution of each signal and prevent cross-signal interference, we apply standardization to each dimension separately before fusing them into a composite reward. This signal is then mapped to token-level advantages for GRPO optimization. By enabling precise, step-wise credit assignment, this design significantly enhances sample efficiency, supporting \texttt{AgentEvolver}'s objective of scalable and continual agent improvement. The entire process is illustrated in the Self-Attributing panel of Figure~\ref{fig:agentevolver_overview}.



\begin{figure}[htbp] 
\centering 


\begin{tcolorbox}[
  title=System Prompt for Step-wise Attribution,
  width=\linewidth, 
  boxrule=0.45mm,
  fonttitle=\bfseries,
  enhanced,
  sharp corners,
  colback=white,
  colframe=black!70,
]
\setlist[itemize]{itemsep=0pt, topsep=2pt, parsep=0pt, partopsep=0pt, leftmargin=*}

You are an expert process reward evaluator specializing in attributional analysis.

\vspace{2mm}
\noindent \textbf{INPUTS:} A user's \texttt{TASK}, a numbered \texttt{SOLUTION TRAJECTORY} of steps, and a final \texttt{OVERALL PERFORMANCE SCORE}.

\vspace{1mm}

\noindent \textbf{YOUR TASK:} Attribute the contribution of each step in the trajectory to the final score based on the following rules.

\vspace{1mm}

\noindent \textbf{EVALUATION RULES:}
\begin{itemize} 
    \item \textbf{If Score is POSITIVE $>0$:}
        \begin{itemize} 
            \item \textbf{GOOD}: The step contributed positively to the solution.
            \item \textbf{BAD}: The step was irrelevant, neutral, or detrimental.
        \end{itemize}
    \item \textbf{If Score is NEGATIVE $\le 0$:}
        \begin{itemize}
            \item \textbf{GOOD}: The step \textbf{only if} it actively corrected or mitigated an error.
            \item \textbf{BAD}: The step introduced, propagated, or failed to fix an error.
        \end{itemize}
\end{itemize}

\vspace{1mm}
\noindent \textbf{FOCUS:} Evaluate strictly on the \textbf{technical impact} of each step's action. Ignore superficial elements like politeness. Reply in the required format only.
\end{tcolorbox}

\caption{System Prompt Structure for step-wise attribution.} 
\label{fig:attr_prompt_system}
\end{figure} 

\subsection{Step-wise Attribution via LLM-based Reasoning} \label{sec:llm_attribution}
The first step of the \texttt{self-attributing} pipeline is to generate the \textbf{process quality} signal (Figure~\ref{fig:agentevolver_overview}, Self-Attributing panel, top). This is accomplished by using an LLM's reasoning capabilities to retrospectively analyze completed trajectories. The goal is to produce a step-wise evaluation of each action's contribution to the final outcome, serving as a measure of its logical correctness within the task context.


To ensure contextual consistency and computational efficiency, we evaluate each trajectory holistically in a single pass. 
Specifically, we format the entire context—including the initial task, all intermediate steps, and the final outcome—into a single prompt (see Figures.~\ref{fig:attr_prompt_system} and~\ref{fig:attr_prompt_user}) for LLM-based evaluation. 
It allows the LLM to capture complex inter-step dependencies, ultimately outputting a binary label for each step: \texttt{GOOD} for beneficial actions and \texttt{BAD} for those deemed irrelevant or counterproductive.








This method of generating the \textbf{process quality} signal offers several advantages over conventional, handcrafted Process Reward Models (PRMs). First, by outputting binary labels (\texttt{GOOD}/\texttt{BAD}), it directly operationalizes the concept of attribution without requiring complex, task-specific scoring schemes, which are often difficult to design and transfer across tasks. Second, it leverages the LLM's high-level reasoning capabilities to make contextual judgments, replacing rigid heuristics with flexible, situation-aware analysis. Most importantly, the resulting attribution labels constitute a semantically-grounded measure of process correctness. They provide directional feedback (i.e., positive or negative contribution) rather than precise reward magnitudes. This qualitative signal is exactly what is needed to construct one of the two independent channels—process quality—before its combination with outcome effectiveness in the subsequent step (Section~\ref{sec:attribution_integration}).
This qualitative signal constitutes the \textbf{process quality} channel, one of two independent inputs to our final reward. Its integration with the 'outcome effectiveness' channel is detailed in Section~\ref{sec:attribution_integration}.


\begin{figure}[htbp] 
\centering 
\begin{tcolorbox}[
  title=User Prompt for Step-wise Attribution,
  width=\linewidth, 
  boxrule=0.45mm,
  fonttitle=\bfseries,
  enhanced,
  sharp corners,
  colback=white,
  colframe=black!70,
]
\textbf{TASK DESCRIPTION} \\
\texttt{<Original task query>}

\vspace{2mm}
\textbf{SOLUTION TRAJECTORY (total N steps)}
\begin{verbatim}
>>> EVAL-STEP 0 <<<
<ACTION> ... <END>
<OBSERVATION> ... <END>
\end{verbatim}
\textit{[…continue for all steps…]}

\vspace{2mm}
\textbf{OVERALL PERFORMANCE SCORE} \texttt{<value>}

\vspace{2mm}
\textbf{REQUIRED OUTPUT FORMAT:}
\begin{description}[style=unboxed, leftmargin=0pt, itemsep=1pt]
    \item[Step 0 Analysis:] \texttt{<your reasoning>}
    \item[Step 0 Judgment:] GOOD/BAD
    
    \textit{[…continue for all steps…]}
\end{description}
\end{tcolorbox}

\caption{User prompt Structure for step-wise attribution, which provides the concrete data for evaluation.} 
\label{fig:attr_prompt_user}
\end{figure}



\subsection{Attribution-Based Reward Construction}
\label{sec:attribution_reward}
Building upon the step-wise qualitative labels (GOOD/BAD) labels generated in Section~\ref{sec:llm_attribution}, this section details the procedure for constructing a quantitative, step-wise attribution reward, $r_t^{\text{attr}}$. 
This signal is designed to function as a dense, process-oriented feedback mechanism that complements the sparse, outcome-based reward. The goal is to convert the LLM's discrete assessment into a continuous signal suitable for policy optimization.

\paragraph{Quantification of Attribution Labels.}

To convert the categorical \texttt{GOOD}/\texttt{BAD} labels into a quantitative signal, we assign a numerical attribution reward, $r_t^{\text{attr}}$, to each step: $+1$ for \texttt{GOOD} and $-1$ for \texttt{BAD}. This binary assignment provides a clear directional signal of each step's localized contribution, without introducing arbitrary scalar magnitudes that could bias gradient computation. Unlike continuous scoring schemes that require careful calibration across diverse tasks and environments, this formulation remains interpretable and stable. These step-wise attribution scores will subsequently be normalized and integrated with outcome-based rewards (Section~\ref{sec:attribution_integration}), enabling the policy to optimize for both procedural correctness and task success.
 

\paragraph{Normalization of the Attribution Signal.} 
To ensure numerical stability and standardize the signal's scale before its integration with the outcome-based reward, we apply standardization to the raw attribution rewards, $r^{\text{attr}}_t$. A critical design choice is the population over which the normalization statistics---mean ($\mu^{\text{attr}}$) and standard deviation ($\sigma^{\text{attr}}$)---are computed. Since trajectories can vary significantly in length, a naive step-level calculation could allow longer, potentially meandering trajectories to disproportionately influence the statistics.

To address this and ensure that each trajectory is treated as an equally important trial, we adopt \textbf{trajectory-level standardization}. In this approach, we first compute the \emph{average} attribution reward for each trajectory. The mean ($\mu^{\text{attr}}$) and standard deviation ($\sigma^{\text{attr}}$) are then calculated over this population of trajectory-average rewards. This method gives each trajectory equal weight in the normalization process, regardless of its length, ensuring that the statistics reflect the distribution of overall trajectory quality rather than being skewed by the volume of individual steps.

The standardized attribution reward, $\hat{r}^{\text{attr}}_t$, is calculated for each step using these trajectory-level statistics:
\begin{equation}
    \hat{r}^{\text{attr}}_t = \frac{r^{\text{attr}}_t - \mu^{\text{attr}}}{\sigma^{\text{attr}} + \epsilon},
\end{equation}
where $\epsilon$ (e.g., $10^{-8}$) is a small constant to prevent division by zero.

Thus, $\hat{r}^{\text{attr}}_t$ serves as a dense, normalized reward signal grounded in the LLM's attribution of each action's correctness.






\subsection{Composite Reward Construction}
\label{sec:attribution_integration}

To form a complete learning signal, we integrate the attribution-based reward $\hat{r}^{\text{attr}}_t$ (from Section~\ref{sec:attribution_reward}) with an outcome-based reward. This dual-channel approach, termed Attribution Outcome Integration (Figure~\ref{fig:agentevolver_overview}, right panel, middle), allows the agent to optimize for both procedural correctness and final task success. The key to this fusion is the seperated normalization of each channel, which preserves their distinct contributions. This section details the formulation of this composite reward.


\paragraph{Outcome Reward Formulation and Normalization.}
To complement the dense, process-oriented attribution signal, we utilize the terminal reward in Section~\ref{sec:self-questioning-llm-judge} as $R^{\text{out}}$, that reflects the final outcome of a trajectory. 

To maintain statistical independence from the attribution channel, this trajectory-level outcome reward is normalized separately across the training group. The standardized outcome reward, $\hat{r}^{\text{out}}$, is computed as:
\begin{equation}
    \hat{r}^{\text{out}} = \frac{R^{\text{out}} - \mu^{\text{out}}}{\sigma^{\text{out}} + \epsilon},
\end{equation}
where $\mu^{\text{out}}$ and $\sigma^{\text{out}}$ are the mean and standard deviation of the raw outcome rewards $R^{\text{out}}$ across all trajectories in the group.
This independent normalization is crucial to prevent the scale of one signal from inadvertently dominating the other during fusion.

\paragraph{Reward Fusion.}
The final composite reward for each step $t$, denoted as $\hat{r}_t$, is then formulated by combining the dense, process-oriented attribution signal with the sparse, outcome-based signal. The outcome reward, $\hat{r}^{\text{out}}$, is applied only at the terminal step of the trajectory to correctly attribute the final success or failure. The formula is a weighted sum:
\begin{equation}
    \hat{r}_t = \alpha \cdot \hat{r}^{\text{attr}}_t + \mathbf{1}_{t=T} \cdot \hat{r}^{\text{out}}.
\end{equation}
Here, $\mathbf{1}_{t=T}$ is an indicator function that equals 1 if $t=T$ and 0 otherwise. The hyperparameter $\alpha \ge 0$ controls the relative contribution of the attribution signal.

\paragraph{Advantage Computation.}
Next, we compute the step-level advantage, $A_t$, from the composite rewards, $\hat{r}_k$. For this, we adopt a simplified yet effective advantage estimation method, following the precedent set by works like DeepSeek-Math~\citep{deepseek_math}. It defines the advantage directly as the undiscounted cumulative future reward. This formulation is a special case of the standard discounted return where the discount factor $\gamma$ is set to 1, thereby giving equal weight to all subsequent steps in a trajectory. The advantage is thus calculated as:
\begin{equation}
    A_t = \sum_{k=t}^{T} \hat{r}_k.
\end{equation}
This calculation produces a single advantage value for each step that cohesively integrates both the process correctness derived from LLM-based attribution and the empirical task success.



\paragraph{The Role of the Weighting Hyperparameter $\alpha$.}
The hyperparameter $\alpha$ provides a direct mechanism to balance the influence of process correctness against outcome maximization. A higher value for $\alpha$ biases the agent towards learning policies that align with the LLM's attribution assessments, fostering robust and generalizable reasoning procedures. Conversely, a lower value prioritizes achieving the final task goal.

This adjustable balance is particularly useful for implementing curriculum learning. For instance, training can commence with a higher $\alpha$ to establish a strong procedural foundation, and then gradually decrease $\alpha$ to fine-tune the policy for optimal task-specific performance.


\subsection{Policy Optimization}
\label{sec:grpo_integration}
The step-level advantages derived in Section~\ref{sec:attribution_integration} must be mapped to token-level signals for policy gradient optimization. This section details the mapping procedure and the resulting optimization objective (Figure~\ref{fig:agentevolver_overview}, right panel, bottom).



\paragraph{Step-to-Token Mapping.}
Policy gradient methods for language models require token-level advantage estimates. For each step $t$ corresponding to an action, we broadcast the step-level advantage $A_t$ to all response tokens generated during that step, yielding token-level advantages $A^{\text{tok}}_j$. Formally, if token $j$ belongs to step $t$ as determined by the step identification mechanism, then $A^{\text{tok}}_j = A_t$.




\paragraph{Integration with Experience-Guided Training.}
These token-level advantages $A^{\text{tok}}_j$ serve as the advantage estimates for both vanilla and experience-guided rollouts in the complete optimization objective (Equation~\ref{eq:navigating_loss}). The composite advantage term incorporates both attribution-based process evaluation and outcome-based task success, enabling the policy to optimize for both intermediate step quality and final trajectory outcomes. This unified treatment ensures consistent credit assignment across different trajectory types while maintaining the selective boosting mechanism for experience-guided samples described in Section~\ref{sec:self-navigating}.

